\begin{abstract}
% video-to-text summarization 보다 더 좋은 표현이 없을까나

Video-to-text summarization remains underexplored in terms of comprehensive evaluation methods. Traditional n-gram overlap-based metrics and recent large language model (LLM)-based approaches depend heavily on human-written reference summaries, limiting their practicality and sensitivity to nuanced semantic aspects. In this paper, we propose QEVA, a reference-free metric evaluating candidate summaries directly against source videos through multimodal question answering. QEVA assesses summaries along three clear dimensions: Coverage, Factuality, and Chronology. We also introduce MLVU(VS)-Eval, a new annotated benchmark derived from the MLVU dataset, comprising 800 summaries generated from 200 videos using state-of-the-art video-language multimodal models. This dataset establishes a transparent and consistent framework for evaluation. Experimental results demonstrate that QEVA shows higher correlation with human judgments compared to existing approaches, as measured by Kendall’s $\tau_b$, $\tau_c$, and Spearman’s $\rho$. We hope that our benchmark and metric will facilitate meaningful progress in video-to-text summarization research and provide valuable insights for the development of future evaluation methods.\footnote{Our evaluation metric (QEVA) is available at \url{https://github.com/jungnerd/QEVA}}

\end{abstract}